% Reviewed source SHA256 (UTF-8/LF): 14cce6c0d191808532ea6ddf8890eaf365d845477df9fdbcc102d7e79791f800
\documentclass[11pt]{article}
\usepackage[T1]{fontenc}
\usepackage{lmodern}
\usepackage[a4paper,margin=27mm]{geometry}
\usepackage{amsmath,amssymb,amsthm,mathtools}
\usepackage{microtype,booktabs,array}
\usepackage[hidelinks]{hyperref}
\usepackage{enumitem}
\setlist{itemsep=3pt,topsep=5pt}
\newtheorem{theorem}{Theorem}[section]
\newtheorem{lemma}[theorem]{Lemma}
\newtheorem{corollary}[theorem]{Corollary}
\newtheorem{proposition}[theorem]{Proposition}
\theoremstyle{definition}
\newtheorem{definition}[theorem]{Definition}
\newtheorem{remark}[theorem]{Remark}
\newcommand{\R}{\mathbb R}
\newcommand{\E}{\mathbb E}
\newcommand{\Prob}{\mathbb P}
\newcommand{\tr}{\operatorname{tr}}
\newcommand{\rank}{\operatorname{rank}}
\newcommand{\diag}{\operatorname{diag}}
\newcommand{\range}{\operatorname{range}}
\newcommand{\OPT}{\operatorname{OPT}}
\newcommand{\norm}[1]{\left\lVert#1\right\rVert}
\newcommand{\ip}[2]{\left\langle#1,#2\right\rangle}
\newcommand{\psd}{\succeq}
\newcommand{\pd}{\succ}
\newcommand{\eps}{\varepsilon}
\DeclareMathOperator*{\argmin}{arg\,min}
\allowdisplaybreaks[1]
\hypersetup{pdfcreator={LaTeX},pdfauthor={Matthew J. Colbrook}}

\title{Convexity of a volume-sampling error sequence}
\author{Matthew J. Colbrook\\\small Department of Applied Mathematics and Theoretical Physics\\\small University of Cambridge, Cambridge, United Kingdom\\\small\href{mailto:m.colbrook@damtp.cam.ac.uk}{m.colbrook@damtp.cam.ac.uk}}
\date{11 September 2026}
\usepackage{xurl}
\setlength{\emergencystretch}{3em}
\begin{document}
\maketitle

\noindent\textbf{Verified affirmative resolution of RA-07.} Theorem 1.1 and its proof; Sections 2--3 establish sampling corollaries.
Independent agent verification is not external human peer review or formal certification. Original draft remarks about review or source access reflect the input date; the dated \href{https://github.com/MColbrook/OpenProblemsInNLA/blob/codex/colbrook-transfer-resolutions/references/colbrook-transfer-2026-09-11/verification/reviews/RA-07-review.md}{independent review} records the subsequent checks.
\par\medskip
% BEGIN REVIEWED BODY
\begin{abstract}
For positive numbers $\lambda_1,\ldots,\lambda_n$, let $e_j$ be their elementary symmetric polynomials, with $e_0=1$ and $e_{n+1}=0$. We prove that $(j+1)e_{j+1}/e_j$ is a decreasing convex sequence on $\{0,\ldots,n\}$. The proof reduces each second difference to the nonnegative quantity $s_1s_3-s_2^2$ for the reciprocal roots of a derivative of the generating polynomial. This proves Conjecture~1 of Derezi\'nski, Khanna, and Mahoney. We give a self-contained application to fixed-size determinantal sampling, remove the concentration loss from its stable-rank bound, and derive a slightly sharper explicit bound.
\end{abstract}

\section{The discrete inequality}
For $0\le j\le n$, write
\[
 e_j=\sum_{\substack{S\subseteq\{1,\ldots,n\}\\|S|=j}}
       \prod_{i\in S}\lambda_i,
 \qquad F_j=(j+1)\frac{e_{j+1}}{e_j}.
\]
Here an empty product is one, and $e_{n+1}=0$. In particular $F_n=0$. The following theorem proves the convexity conjecture stated in \cite[Section~5, Conjecture~1]{DKM2020}.

\begin{theorem}\label{thm:convex}
If every $\lambda_i$ is positive, then $F_0>F_1>\cdots>F_n$ and
\[
 F_{k-1}-2F_k+F_{k+1}\ge0\qquad(1\le k\le n-1).
\]
More precisely, factor
\[
 Q(t):=\frac{d^{k-1}}{dt^{k-1}}\prod_{i=1}^n(1+\lambda_i t)
     =Q(0)\prod_{a=1}^{n-k+1}(1+\mu_a t),
 \qquad s_\ell=\sum_a\mu_a^\ell.
\]
All $\mu_a$ are positive, and
\begin{equation}\label{eq:certificate}
 F_{k-1}-2F_k+F_{k+1}
 =\frac{2(s_1s_3-s_2^2)}{s_1(s_1^2-s_2)}
 =\frac{2\sum_{a<b}\mu_a\mu_b(\mu_a-\mu_b)^2}
        {s_1(s_1^2-s_2)}.
\end{equation}
\end{theorem}

\begin{proof}
Put $P(t)=\prod_i(1+\lambda_i t)$. Since $P^{(j)}(0)=j!e_j$,
\[
 F_j=\frac{P^{(j+1)}(0)}{P^{(j)}(0)}.
\]
Every root of $P$ is strictly negative. Repeated applications of Rolle's theorem, including the usual multiplicities at repeated roots, show that every derivative of positive degree also has only strictly negative roots. Thus $Q$ has the stated factorization with positive $\mu_a$.

The first three derivatives of this product at zero give
\[
 \frac{Q'(0)}{Q(0)}=s_1,\qquad
 \frac{Q''(0)}{Q(0)}=s_1^2-s_2,\qquad
 \frac{Q'''(0)}{Q(0)}=s_1^3-3s_1s_2+2s_3.
\]
Consequently
\[
 F_{k-1}=s_1,\qquad
 F_k=\frac{s_1^2-s_2}{s_1},\qquad
 F_{k+1}=\frac{s_1^3-3s_1s_2+2s_3}{s_1^2-s_2}.
\]
Subtracting gives the first identity in \eqref{eq:certificate}. Its denominator is positive because $Q$ has degree at least two and
$s_1^2-s_2=2\sum_{a<b}\mu_a\mu_b>0$. The numerator satisfies
\[
 s_1s_3-s_2^2=\sum_{a<b}\mu_a\mu_b(\mu_a-\mu_b)^2\ge0.
\]
This proves convexity, including $k=n-1$: in that case $Q$ has degree two and $Q'''=0$.

For monotonicity, use the same calculation with $Q=P^{(j)}$, for $0\le j\le n-1$. It gives $F_j-F_{j+1}=s_2/s_1>0$. When $Q$ has degree one, this identity still holds, with $F_{j+1}=F_n=0$.
\end{proof}

If all $\lambda_i$ equal $\lambda$, then $F_j=(n-j)\lambda$ and every second difference vanishes. Formula~\eqref{eq:certificate} also supplies an explicit nonnegative certificate in the general case; no limiting or probabilistic argument is needed for the theorem.

\section{Determinantal sampling}
Let $A\in\R^{m\times N}$ have rank $r$, and let $\lambda_1\ge\cdots\ge\lambda_r>0$ be the nonzero eigenvalues of $G=A^TA$. Define
\[
 \mathcal E(S)=\norm{A-P_SA}_F^2,
\]
where $P_S$ is the orthogonal projector onto the span of the columns indexed by $S$. For $0\le j\le r$, the fixed-size determinantal distribution assigns probability proportional to $\det(G_{S,S})$ to subsets with $|S|=j$.

\begin{lemma}\label{lem:expectations}
For fixed-size sampling of size $j$,
\[
 \E_j\mathcal E(S)=F_j=(j+1)e_{j+1}/e_j.
\]
For the random-size distribution
\[
 \Prob_\alpha(S)=\frac{\alpha^{-|S|}\det(G_{S,S})}{\det(I+G/\alpha)},
 \qquad \alpha>0,
\]
write $K=|S|$. Then
\begin{equation}\label{eq:dpp}
 \mu(\alpha):=\E_\alpha K=\sum_{i=1}^r\frac{\lambda_i}{\lambda_i+\alpha},
 \qquad \E_\alpha\mathcal E(S)=\alpha\mu(\alpha).
\end{equation}
\end{lemma}

\begin{proof}
The principal-minor identity gives
$\sum_{|S|=j}\det(G_{S,S})=e_j$.
For $i\notin S$, the Gram determinant formula gives
\[
 \det(G_{S,S})\norm{(I-P_S)a_i}_2^2
       =\det(G_{S\cup\{i\},S\cup\{i\}}).
\]
This also holds when $G_{S,S}$ is singular: adding one column cannot increase its rank to $j+1$, so both sides vanish. Summing over $S$ and $i$ counts every $(j+1)$-set exactly $j+1$ times, proving the fixed-size formula.

Put $P(t)=\sum_{j=0}^r e_jt^j=\prod_{i=1}^r(1+\lambda_i t)$. The probability generating function of $K$ has normalizing factor $P(1/\alpha)$. Thus
\[
 \E_\alpha K=\frac{\alpha^{-1}P'(1/\alpha)}{P(1/\alpha)},\qquad
 \E_\alpha\mathcal E(S)=\frac{P'(1/\alpha)}{P(1/\alpha)}.
\]
The logarithmic derivative of $P$ proves \eqref{eq:dpp}.
\end{proof}

\begin{corollary}\label{cor:jensen}
Let $1\le k<r$. If $\mu(\alpha)\le k$, then
\[
 \E_k\mathcal E(S)\le\alpha\mu(\alpha)\le\alpha k.
\]
In particular, there is a unique $\alpha_k>0$ satisfying $\mu(\alpha_k)=k$, and
\[
 \E_k\mathcal E(S)\le\E_{\alpha_k}\mathcal E(S)=\alpha_k k.
\]
\end{corollary}

\begin{proof}
Let $F(x)$ be the piecewise-linear interpolation of the values $F_j$ on $[0,r]$. Theorem~\ref{thm:convex} makes $F$ convex and decreasing. Conditional on $K=j$, random-size sampling is the fixed-size distribution, so
\[
 F(k)\le F(\E K)\le\E F(K)=\E_\alpha\mathcal E(S).
\]
The remaining assertions follow from \eqref{eq:dpp} and the strict decrease of $\mu(\alpha)$ from $r$ to zero as $\alpha$ increases from zero to infinity.
\end{proof}

\section{Stable-rank bounds without a concentration loss}
For $0\le s<r$, define the order-$s$ stable rank by
\[
 r_s=\frac{\sum_{i>s}\lambda_i}{\lambda_{s+1}},\qquad t_s=s+r_s.
\]
The optimal rank-$k$ squared Frobenius error is
$T_k=\sum_{i>k}\lambda_i$.

\begin{theorem}\label{thm:stable}
For every integer $s<k<t_s$, fixed-size determinantal sampling satisfies
\begin{align}
 \frac{\E_k\mathcal E(S)}{T_k}
 &\le \frac{k}{2(k-s)}
      \left(1+\sqrt{1+\frac{4(k-s)\lambda_{s+1}}{T_k}}\right)\label{eq:directbound}\\
 &\le \frac{k}{2(k-s)}
      \left(1+\sqrt{1+\frac{4(k-s)}{t_s-k}}\right)\label{eq:stablebound}\\
 &\le \frac{k}{k-s}\sqrt{1+\frac{2(k-s)}{t_s-k}}.\label{eq:oldphi}
\end{align}
No auxiliary $\eps$, lower bound on $k-s$, or concentration-loss multiplier is required.
\end{theorem}

\begin{proof}
Write $\ell=k-s$, $\beta=\lambda_{s+1}$, and $T=T_k$. Since the eigenvalues are decreasing,
\[
 \mu(\alpha)\le s+\frac{\ell\beta}{\beta+\alpha}+\frac{T}{\alpha}.
\]
Choose
\[
 \alpha=\frac{T+\sqrt{T^2+4\ell\beta T}}{2\ell}.
\]
It solves $\ell\alpha^2=T(\alpha+\beta)$, so the preceding upper bound on $\mu(\alpha)$ equals $s+\ell=k$. Corollary~\ref{cor:jensen} gives $F_k\le k\alpha$, which is \eqref{eq:directbound}.

Moreover,
\[
 T=\beta r_s-\sum_{i=s+1}^k\lambda_i
       \ge\beta(r_s-\ell)=\beta(t_s-k)>0.
\]
This proves \eqref{eq:stablebound}. Finally,
\[
 \frac{1+\sqrt{1+4x}}2\le\sqrt{1+2x}\qquad(x\ge0)
\]
follows by squaring and using $\sqrt{1+4x}\le1+2x$. Substituting $x=(k-s)/(t_s-k)$ proves \eqref{eq:oldphi}.
\end{proof}

The final expression in \eqref{eq:oldphi} is the function $\Phi_s(k)$ appearing in \cite[Theorem~1]{DKM2020}. Theorem~\ref{thm:stable} therefore yields the loss-free version contemplated there, over the full range $s<k<t_s$, with the sharper intermediate bounds \eqref{eq:directbound} and \eqref{eq:stablebound}.

For completeness, the classical worst-case bound $F_k\le(k+1)T_k$ follows directly from $e_{k+1}\le T_ke_k$: every $(k+1)$-set contains an index greater than $k$, and expanding $T_ke_k$ counts its product at least once. Thus one may take the minimum of $k+1$ and any applicable bound above.

\section{Scope and verification}
The polynomial theorem concerns positive numbers; a rank-deficient matrix is handled by using only its positive eigenvalues, as in Section~2. Fixed-size determinantal sampling is asserted only for sizes at most the rank. The endpoint $k=r$ has zero residual and needs no relative-error ratio.

The accompanying script \texttt{verification/verify\_volume\_sampling.py} checks the second-difference identity symbolically, checks the polynomial inequalities in exact rational arithmetic, and checks the determinantal expectation identities on small matrices. These checks supplement, rather than replace, the proofs above.

\begin{thebibliography}{9}
\bibitem{DKM2020}
M. Derezi\'nski, R. Khanna, and M. W. Mahoney.
\newblock Improved guarantees and a multiple-descent curve for Column Subset Selection and the Nystr\"om method.
\newblock \emph{arXiv:2002.09073}, version~3, 18 December 2020.
\newblock \url{https://arxiv.org/abs/2002.09073}.
\end{thebibliography}
\end{document}

% END REVIEWED BODY
